% ======================================================================
% 第 3 章：格点QCD计算的范围
% 原著：R. Gupta, "Introduction to Lattice QCD" (arXiv:hep-lat/9807028)，
% 第 3 节，第 10–14 页。结尾段落与项目符号列表延续至第 14 页
% （其文字摘录于 sec04.txt 顶部）。
% ======================================================================

\chapter{格点QCD计算的范围}

\begin{figure}[htbp]
\centering
\includegraphics[width=0.8\textwidth]{images/fig1.png}
\caption{Wolfenstein 参数化下的 CKM 矩阵。我展示了一些物理过程的例子，这些过程
需要在实验上测量，并通过「主」方程与理论预言进行比较，以估计 CKM 矩阵的各个
矩阵元。}
\label{fig:1}
\end{figure}

LQCD 的首要目标是检验 QCD 是否是强相互作用的正确理论。大多数粒子物理学家相信答案是
肯定的，因为它在描述大动量转移过程方面取得了成功，并且具有数学上的优美性和唯一性。
然而，定量的确认仍然缺乏。在分析尺度 $\lesssim 1$~GeV 下的强子过程时，传统的微扰方法
（按小参数展开）的问题在于强耦合常数 $\alpha_s \sim 1$。因此，按 $\alpha_s$ 的
微扰论并不可靠。结果是，即使给出 $\alpha_s$ 和夸克的质量，我们也无法从 QCD 计算
介子和重子的质量。图~\ref{fig:2} 对此作了说明，我在此比较了氢原子的结合能与质子的
结合能。结果表明，质子的质量几乎全部归因于胶子的强非线性相互作用。

\begin{figure}[htbp]
\centering
\includegraphics[width=0.8\textwidth]{images/fig2.png}
\caption{氢原子中电子与质子之间（通过电磁力相互作用）的结合能，与质子中夸克
（强力）的结合能之差异的示意图。}
\label{fig:2}
\end{figure}

格点QCD 的模拟有六个未知的输入参数。它们是耦合常数 $\alpha_s$ 以及上、下、奇异、
粲和底夸克的质量（顶夸克的寿命太短，$0.4 \times 10^{-24}$ 秒，不足以形成束缚态，
因此不适合用于格点研究）。一旦用六个精确测量的强子质量把它们确定下来，此后由这些
夸克和胶子构成的数百种其他粒子的质量与性质都必须与实验相符。例如，考虑 $\pi$ 介子质量
的手征微扰论（$\chi$PT）关系
\begin{equation}
(M_{\pi^+})^2 = B_\pi (m_u + m_d) + C_\pi (m_u + m_d)^2 + \ldots
\label{eq:chipt}
\end{equation}
格点计算将为这一关系提供检验，并确定常数 $B_\pi$、$C_\pi$、……，即约束手征拉氏量。
此外，在测量了一组夸克质量值下的 $M_\pi$ 之后，可以对式~\eqref{eq:chipt} 求逆以确定
流夸克质量。我把对 QCD 的这种检验以及对 SM 中六个未知参数的测定，称为格点QCD 的
直接应用。

微扰估计不可靠的第二个领域，是强子弱衰变中所出现矩阵元的计算。其原因在于，对基本
弱过程而言，由于初态与末态之间软胶子的交换，非微扰 QCD 修正可能很大。图~\ref{fig:3}
以半轻子衰变 $D \to K l\nu$ 为例对此作了说明。在这种情况下，初态 $D$ 介子与末态
K 介子之间的弱相互作用哈密顿量矩阵元，会收到来自 QCD 的大修正，而这些修正无法用
PQCD 来估计。

\begin{figure}[htbp]
\centering
\includegraphics[width=0.8\textwidth]{images/fig3.png}
\caption{半轻子衰变 $D^- \to K^0 e^- \bar{\nu}_e$ 的费曼图。QCD 修正由初态与
末态强子之间交换的各种胶子来说明。轻子顶点可以用微扰论可靠地计算，而强子顶点
则需要非微扰方法。}
\label{fig:3}
\end{figure}

实验的结果是一个数，即衰变率。为了检验标准模型的预言是否与此数相符，理论家必须为
同一过程推导出解析表达式。这一表达式通常由三部分的乘积组成：SM 中参数的组合、相互
作用哈密顿量在初态与末态之间的矩阵元（ME），以及运动学因子。（这就是著名的费米黄金
定则。）示意性地，可以把这一「主」方程写成如下形式。
\begin{equation}
\mathrm{Expt.}\# =
(\mathrm{SM\ parameters})\,(\mathrm{matrix\ elements})\,(\mathrm{kinematic\ factors}).
\label{eq:master}
\end{equation}
因此，对于每一个这样的实验数值，只要 ME 已知，就能得到对 SM 参数某一特定组合的
一个约束。

通过利用若干精密实验并计算必要的 ME，可以对模型进行过度约束，从而检验它。我把对 ME
的格点计算称为对标准模型参数的「间接」探测，因为人们通常需要额外的现象学/实验输入，
才能通过主方程提取未知参数。

过去 25 年来，物理学家们一直在忙于计算和测量这些参数，因此人们已经有了一些合理的
估计。因此，记住格点计算若要产生影响所必须达到的精度水平是很有用的。非格点估计——
取自 1996 年《粒子性质评论》(Review of Particle Properties) 以及 Prof.~Buras 的
讲义 [10]——如下

\begin{table}[htbp]
\centering
\caption*{标准模型参数的非格点估计}
\label{tab:nl-estimates}
\begin{tabular}{l l p{6.5cm}}
\toprule
参数 & 数值 & 说明 \\
\midrule
$m_u(\overline{MS}, 2\ \mathrm{GeV})$ & 3--4.5~MeV & $\chi$PT，求和规则 \\
$m_d(\overline{MS}, 2\ \mathrm{GeV})$ & 5--8~MeV & $\chi$PT，求和规则 \\
$m_s(\overline{MS}, 2\ \mathrm{GeV})$ & 100--140~MeV & $\chi$PT，求和规则 \\
$m_c(\overline{MS}, m_c)$ & 1.0--1.6~GeV & $J/\psi$ 谱 \\
$m_b(\overline{MS}, m_b)$ & 4.1--4.6~GeV & $\Upsilon$ 谱 \\
$M_{\mathrm{Higgs}}$ & & 未找到 \\
$\alpha_s(M_Z)$ & $0.118 \pm 0.003$ & 世界平均值 \\
$|V_{ud}|$ & $0.9736 \pm 0.0010$ & \begin{tabular}[t]{@{}l@{}}
$n \to p e\nu_e$ 和 \\ $\mu \to e\nu_e\nu_\mu$ \end{tabular} \\
$|V_{us}|$ & $0.2196 \pm 0.0023$ & \begin{tabular}[t]{@{}l@{}}
$K^+ \to \pi^0 e^+ \nu_e$ 和 \\ $K_L^0 \to \pi^\pm e^\mp \nu_e$ \end{tabular} \\
$|V_{ub}|$ & $(3.3 \pm 0.2 \pm 0.4 \pm 0.7) \times 10^{-3}$ &
\begin{tabular}[t]{@{}l@{}} $B \to \pi l\nu$; $B \to \rho l\nu$ \end{tabular} \\
$|V_{cb}|$ & $0.040 \pm 0.003$ & \begin{tabular}[t]{@{}l@{}}
$B \to X_c l\nu$ 和 \\ $B \to D^* l\nu$ \end{tabular} \\
$|\epsilon|$ & $2.26 \times 10^{-3}$ & $K^0 \leftrightarrow \bar{K}^0$ 混合 \\
$|\epsilon/\epsilon'|$ & $0$--$30 \times 10^{-4}$ & $K \to \pi\pi$ 衰变 \\
\bottomrule
\end{tabular}
\end{table}

注意，实验只能确定 CKM 矩阵元的模。显然，格点QCD 最具吸引力的机遇是帮助确定夸克
质量、$\alpha_s$ 以及基本上未知的 Wolfenstein 参数 $\rho$ 和 $\eta$。对于 $V_{CKM}$
的其他矩阵元，人们已有合理的估计，然而在这里，进一步的进展也需要强子矩阵元估计的
显著改进。现象学家正依赖 LQCD 来提供改进的估计。

让我以如下内容结束本节：除了通过再现所测量的介子、重子和胶球谱以及计算弱矩阵元来
验证 QCD 是强相互作用的正确理论之外，LQCD 的模拟还将使我们能够
\begin{itemize}
\item 研究 QCD 真空的拓扑结构，以及负责禁闭和自发破缺手征对称性的机制。近期综述见 [12]。
\item 计算诸如波函数、衰变常数、电磁、弱和稀有衰变的形状因子，以及深度非弹性散射的
结构函数等强子性质。通过计算矩阵元 $\langle N | q\bar{q} | N \rangle$ 和
$\langle N | q\gamma_\mu \gamma_5 q | N \rangle$，估计 $\pi$ 介子-核子 $\sigma$ 项以及夸克所
携带的质子自旋。其中一些课题将由 Prof.~Martinelli 讲解，其余的课题一个好的参考是
LATTICE 会议文集 [9]。
\item 分析 QCD 在高温下的行为，并处理诸如相变的性质、相变温度 $T_c$、状态方程，
以及 $T_c$ 附近及其以上的集体激发等问题。本学校 J-P Blaizot 已讲解了有限温度 QCD
这一课题，因此我将只为你列出我偏爱的关于该主题的教学性综述 [13]。Profs.~Martinelli、
L\"{u}scher 和我的讲义将仅限于理解和探测零温下 QCD 的性质。
\end{itemize}
